\begin{explanationbox}
\begin{center}
\begin{minipage}{0.98\linewidth}
\centering
\textbf{一项增，一项减：和在相等处最低}\par\medskip
\begin{tikzpicture}[font=\small,>=Stealth]

\begin{scope}[x=2.1cm,y=0.95cm]
\draw[->] (0,0)--(3.65,0) node[right] {$x$};
\draw[->] (0,0)--(0,4.5) node[above] {数值};
\draw[CStatement,very thick,domain=0.26:3.3,samples=90] plot (\x,{\x+1/\x});
\draw[CProof,thick,dashed,domain=0.26:3.3] plot (\x,\x);
\draw[CExample,thick,dotted,domain=0.26:3.3,samples=90] plot (\x,{1/\x});
\draw[gray,dashed] (1,0)--(1,2);
\fill[CProof] (1,1) circle (2.5pt); \fill[CStatement] (1,2) circle (2.5pt);
\node[below] at (1,0) {$1$};
\node[right] at (2.6,3.85) {$x+1/x$};
\node[right] at (3.1,2.85) {$x$};
\node[right] at (3.1,0.55) {$1/x$};
\node[above right,fill=white] at (1,2) {和为 $2$};
\node[below right,fill=white] at (1.1,1.15) {两项均为 $1$};
\end{scope}
\end{tikzpicture}\par\smallskip
先看 $x+1/x$：两项的乘积始终为 $1$。推广到 $Ax$ 与 $B/x$，乘积为 $AB$，相等处的和为 $2\sqrt{AB}$。
\end{minipage}
\end{center}


图中 $x$ 增大时，$x$ 增大而 $1/x$ 减小，乘积保持为 $1$；两项相等时，和最低。
一般情形的取等条件为 $Ax=B/x$，正数解为 $x=\sqrt{B/A}$。
负半轴利用 $f(-x)=-f(x)$，把正半轴的最低点映到负半轴的最高点。
求值域前先检查 $A,B>0$，以及定义域是否允许取等；分式换元后也要重新检查新变量的范围。

\medskip
\textbf{\textcolor{CExplanation}{几何迁移：由对勾最值到两支曲线的最短距离}}

\textbf{题目：} $k>0$，点 $P$ 在 $y=\frac{k}{x}$ 第一象限部分，点 $Q$ 在第三象限部分，求 $PQ$ 的最小值。

\begin{center}
\begin{minipage}{0.98\linewidth}
\centering
\textbf{先分清水平差、竖直差与真实距离}\par\medskip
\begin{tikzpicture}[font=\small,>=Stealth]

\begin{scope}[x=1.65cm,y=1.65cm]
\draw[->,gray] (-1.7,0)--(2.7,0) node[right] {$x$};
\draw[->,gray] (0,-1.6)--(0,1.7) node[above] {$y$};
\draw[CProof,very thick] (-1,-1)--(2,0.5);
\draw[CExample,thick] (-1,-1)--(2,-1)--(2,0.5);
\draw (1.85,-1)--(1.85,-0.85)--(2,-0.85);
\fill (-1,-1) circle (1.7pt) (2,0.5) circle (1.7pt);
\node[above right] at (2,0.5) {$P(a,k/a)$};
\node[below left] at (-1,-1) {$Q(-b,-k/b)$};
\node[below] at (0.5,-1) {$u=a+b$};
\node[right] at (2,-0.45) {$v=k/a+k/b$};
\node[above,rotate=26.565] at (0.6,0.25) {$d=PQ$};
\end{scope}
\end{tikzpicture}\par\smallskip
图取 $k=1,a=2,b=1$。两条直角边为坐标差，斜边才是距离；$d^2\ge2uv$ 的取等要求是 $u=v$。
\end{minipage}
\end{center}

设 $P(a,\frac{k}{a})$、$Q(-b,-\frac{k}{b})$，其中 $a>0$、$b>0$。这样取 $b>0$，可避免第三象限坐标的符号混乱。
\[
\Delta x=a+b>0,\qquad
\Delta y=\frac{k}{a}+\frac{k}{b}>0.
\]
\textbf{关键观察：} 与其直接求最小距离，不如先看两坐标差之和：
\[
\Delta x+\Delta y
=\underbrace{\left(a+\frac{k}{a}\right)}_{\ge2\sqrt{k}}
+\underbrace{\left(b+\frac{k}{b}\right)}_{\ge2\sqrt{k}}
\ge4\sqrt{k}.
\]
因为 $2(\Delta x^2+\Delta y^2)-(\Delta x+\Delta y)^2=(\Delta x-\Delta y)^2\ge0$，所以
\[
PQ=\sqrt{\Delta x^2+\Delta y^2}
\ge\frac{\Delta x+\Delta y}{\sqrt2}
\ge2\sqrt{2k}.
\]
\textbf{几何意义：} $\frac{\Delta x+\Delta y}{\sqrt2}$ 是连接线段沿 $y=x$ 方向走过的长度；只有线段本身也平行 $y=x$，这部分长度才等于整段线段的长度。

\textbf{两次等号同时成立：} $a+\frac{k}{a}=2\sqrt{k}$ 与 $b+\frac{k}{b}=2\sqrt{k}$ 要求 $a=b=\sqrt{k}$；此时 $\Delta x=\Delta y=2\sqrt{k}$，所以确实取得 $PQ=2\sqrt{2k}$.

\textbf{识别信号：} 两点分处第一、第三象限；函数满足 $xy=k>0$；题目求两支间的\textbf{欧氏距离}。先想到 $x+y=x+\frac{k}{x}$，再结合勾股定理；切忌把某一点到原点的最小距离直接当作两支间距离的一半，而不说明理由。

\end{explanationbox}
